\documentclass[10pt,a4paper]{amsart}
\usepackage{amsmath,amssymb,amsthm,graphicx}
\usepackage{verbatim}
\newenvironment{reference}{\comment}{\endcomment}
%\newenvironment{reference}{}{}
\newenvironment{slogan}{\comment}{\endcomment}
\newenvironment{history}{\comment}{\endcomment}
% For commutative diagrams we use Xy-pic
\usepackage[all]{xy}

% We use 2cell for 2-commutative diagrams.
\xyoption{2cell}
% Theorem environments.
%
\theoremstyle{plain}
\newtheorem{theorem}[subsection]{Theorem}
\newtheorem{proposition}[subsection]{Proposition}
\newtheorem{lemma}[subsection]{Lemma}

\theoremstyle{definition}
\newtheorem{definition}[subsection]{Definition}
\newtheorem{example}[subsection]{Example}
\newtheorem{exercise}[subsection]{Exercise}
\newtheorem{situation}[subsection]{Situation}

\theoremstyle{remark}
\newtheorem{remark}[subsection]{Remark}
\newtheorem{remarks}[subsection]{Remarks}

\numberwithin{equation}{subsection}

% Macros
%
\def\lim{\mathop{\mathrm{lim}}\nolimits}
\def\colim{\mathop{\mathrm{colim}}\nolimits}
\def\Spec{\mathop{\mathrm{Spec}}}
\def\Hom{\mathop{\mathrm{Hom}}\nolimits}
\def\Ext{\mathop{\mathrm{Ext}}\nolimits}
\def\SheafHom{\mathop{\mathcal{H}\!\mathit{om}}\nolimits}
\def\SheafExt{\mathop{\mathcal{E}\!\mathit{xt}}\nolimits}
\def\Sch{\mathit{Sch}}
\def\Mor{\mathop{\mathrm{Mor}}\nolimits}
\def\Ob{\mathop{\mathrm{Ob}}\nolimits}
\def\Sh{\mathop{\mathit{Sh}}\nolimits}
\def\NL{\mathop{N\!L}\nolimits}
\def\CH{\mathop{\mathrm{CH}}\nolimits}
\def\proetale{{pro\text{-}\acute{e}tale}}
\def\etale{{\acute{e}tale}}
\def\QCoh{\mathit{QCoh}}
\def\Ker{\mathop{\mathrm{Ker}}}
\def\Im{\mathop{\mathrm{Im}}}
\def\Coker{\mathop{\mathrm{Coker}}}
\def\Coim{\mathop{\mathrm{Coim}}}

% Boxtimes
%
\DeclareMathSymbol{\boxtimes}{\mathbin}{AMSa}{"02}

%
% Macros for moduli stacks/spaces
%
\def\QCohstack{\mathcal{QC}\!\mathit{oh}}
\def\Cohstack{\mathcal{C}\!\mathit{oh}}
\def\Spacesstack{\mathcal{S}\!\mathit{paces}}
\def\Quotfunctor{\mathrm{Quot}}
\def\Hilbfunctor{\mathrm{Hilb}}
\def\Curvesstack{\mathcal{C}\!\mathit{urves}}
\def\Polarizedstack{\mathcal{P}\!\mathit{olarized}}
\def\Complexesstack{\mathcal{C}\!\mathit{omplexes}}
% \Pic is the operator that assigns to X its picard group, usage \Pic(X)
% \Picardstack_{X/B} denotes the Picard stack of X over B
% \Picardfunctor_{X/B} denotes the Picard functor of X over B
\def\Pic{\mathop{\mathrm{Pic}}\nolimits}
\def\Picardstack{\mathcal{P}\!\mathit{ic}}
\def\Picardfunctor{\mathrm{Pic}}
\def\Deformationcategory{\mathcal{D}\!\mathit{ef}}
\usepackage{cleveref}
\crefname{originallemma}{original lemma}{original lemmas}
\Crefname{originallemma}{Original lemma}{Original lemmas}
\makeatletter
\newlabel{lemma-inertia}{{lemma-inertia}{source}{}{sourcepointer-1}{}}
\newlabel{lemma-inertia@cref}{{[originallemma][0][]lemma-inertia}{[1][0][]source}{}{}{}}
\makeatother
\crefname{originallemma}{original lemma}{original lemmas}
\Crefname{originallemma}{Original lemma}{Original lemmas}
\makeatletter
\newlabel{lemma-isom-pseudo-torsor-aut-over-space}{{lemma-isom-pseudo-torsor-aut-over-space}{source}{}{sourcepointer-2}{}}
\newlabel{lemma-isom-pseudo-torsor-aut-over-space@cref}{{[originallemma][0][]lemma-isom-pseudo-torsor-aut-over-space}{[1][0][]source}{}{}{}}
\makeatother
\crefname{originalremark}{original remark}{original remarks}
\Crefname{originalremark}{Original remark}{Original remarks}
\makeatletter
\newlabel{remark-inertia-is-group-in-spaces}{{remark-inertia-is-group-in-spaces}{source}{}{sourcepointer-3}{}}
\newlabel{remark-inertia-is-group-in-spaces@cref}{{[originalremark][0][]remark-inertia-is-group-in-spaces}{[1][0][]source}{}{}{}}
\makeatother
\crefname{originalreference}{original reference}{original references}
\Crefname{originalreference}{Original reference}{Original references}
\makeatletter
\newlabel{spaces-groupoids-lemma-group-scheme-separated}{{spaces-groupoids-lemma-group-scheme-separated}{source}{}{sourcepointer-4}{}}
\newlabel{spaces-groupoids-lemma-group-scheme-separated@cref}{{[originalreference][0][]spaces-groupoids-lemma-group-scheme-separated}{[1][0][]source}{}{}{}}
\makeatother
\crefname{originalreference}{original reference}{original references}
\Crefname{originalreference}{Original reference}{Original references}
\makeatletter
\newlabel{spaces-groupoids-lemma-pseudo-torsor-implications}{{spaces-groupoids-lemma-pseudo-torsor-implications}{source}{}{sourcepointer-5}{}}
\newlabel{spaces-groupoids-lemma-pseudo-torsor-implications@cref}{{[originalreference][0][]spaces-groupoids-lemma-pseudo-torsor-implications}{[1][0][]source}{}{}{}}
\makeatother
\title{Separatedness, unramifiedness and local quasi-finiteness of the diagonal read off the relative inertia}
\author{The Stacks Project Authors}
\date{stacks-project:stacks-morphisms.tex:local-snapshot-20261012; GNU FDL 1.2 or later}
\begin{document}
\maketitle
\noindent\textit{Editorial note.} All delivered mathematical text is the author's own, copied from the pinned Stacks Project chapter source sources/stacks-algebraic-stacks-morphisms/pinned-original-source.tex; no statement or proof step was rewritten, reordered or removed.\par\medskip
\noindent\textit{Editorial note.} Preamble provenance: the Stacks Project keeps one single shared preamble for every chapter; this card takes it directly from the staged sources/stacks-algebraic-stacks-morphisms/preamble.tex as three verbatim spans (lines 8-12, 14-18 and 154-217), so no derived file was written and nothing was added or redefined. The wrapper class and the base packages are supplied by the assembler.\par\medskip
\noindent\textit{Editorial note.} Bibliography provenance: neither the counted statement nor its proof contains a citation (this lane rescans the whole span before writing the card), so the card carries no bibliography and no entry was reconstructed.\par\medskip
\noindent\textit{Editorial note.} Selection provenance: this card was opened under the file-based control contract (CONTROL.md), with the mechanical pre-card gate run on its statement and proof interval before the card was written.\par\medskip
\noindent\textit{Allowed standard prerequisites (editorial).}\begin{itemize}
\item The construction of the relative inertia as a group algebraic space over a presentation of the source. Invoked inside the counted proof (source lines 1642-1647). Turns the identity section of the relative inertia into the identity of a group space.
\item The statement that a group scheme over a base is separated exactly when its identity section is a closed immersion. Invoked inside the counted proof (source lines 1648-1651) as a cross-chapter pointer. Gives the equivalence of the three conditions in part (1).
\item The description of the diagonal through the isomorphism spaces of a presentation. Used inside the counted proof (source lines 1655-1685). Identifies the pullbacks of the two diagonals with isomorphism spaces.
\item The statement that the isomorphism space of a pseudo torsor is unramified or locally quasi-finite when the group space is. Invoked at the end of the counted proof (source lines 1686-1695) as a cross-chapter pointer. Transfers the properties from the inertia to the diagonal in parts (4) and (5).
\end{itemize}\medskip
\ifdefined\hypertarget\hypertarget{sourcepointer-1}{}\fi
\noindent\textit{Source reference.} Original lemma lemma-inertia: sources/stacks-algebraic-stacks-morphisms/pinned-original-source.tex, source lines 1220--1220. Navigation pointer only: the counted text refers to the Stacks item lemma-inertia, which is an existing labeled result of this same chapter and is not part of this card. The pointer carries the target's own label as its number and its definition line as its span; its content is not reproduced here.\par\medskip
\ifdefined\hypertarget\hypertarget{sourcepointer-2}{}\fi
\noindent\textit{Source reference.} Original lemma lemma-isom-pseudo-torsor-aut-over-space: sources/stacks-algebraic-stacks-morphisms/pinned-original-source.tex, source lines 1392--1392. Navigation pointer only: the counted text refers to the Stacks item lemma-isom-pseudo-torsor-aut-over-space, which is an existing labeled result of this same chapter and is not part of this card. The pointer carries the target's own label as its number and its definition line as its span; its content is not reproduced here.\par\medskip
\ifdefined\hypertarget\hypertarget{sourcepointer-3}{}\fi
\noindent\textit{Source reference.} Original remark remark-inertia-is-group-in-spaces: sources/stacks-algebraic-stacks-morphisms/pinned-original-source.tex, source lines 1292--1292. Navigation pointer only: the counted text refers to the Stacks item remark-inertia-is-group-in-spaces, which is an existing labeled result of this same chapter and is not part of this card. The pointer carries the target's own label as its number and its definition line as its span; its content is not reproduced here.\par\medskip
\ifdefined\hypertarget\hypertarget{sourcepointer-4}{}\fi
\noindent\textit{Source reference.} Original reference spaces-groupoids-lemma-group-scheme-separated: sources/stacks-algebraic-stacks-morphisms/pinned-original-source.tex, source lines 1658--1658. Navigation pointer only: the counted text refers to the Stacks item spaces-groupoids-lemma-group-scheme-separated, which belongs to another chapter of the project and is not part of this card. The pointer carries the target's own label as its number because this chapter's staged snapshot does not contain the other chapter's tag table; its content is not reproduced here.\par\medskip
\ifdefined\hypertarget\hypertarget{sourcepointer-5}{}\fi
\noindent\textit{Source reference.} Original reference spaces-groupoids-lemma-pseudo-torsor-implications: sources/stacks-algebraic-stacks-morphisms/pinned-original-source.tex, source lines 1697--1697. Navigation pointer only: the counted text refers to the Stacks item spaces-groupoids-lemma-pseudo-torsor-implications, which belongs to another chapter of the project and is not part of this card. The pointer carries the target's own label as its number because this chapter's staged snapshot does not contain the other chapter's tag table; its content is not reproduced here.\par\medskip
\begin{lemma}
\label{lemma-diagonal-diagonal}
Let $f : \mathcal{X} \to \mathcal{Y}$ be a morphism of algebraic stacks.
\begin{enumerate}
\item
The following are equivalent
\begin{enumerate}
\item $\mathcal{I}_{\mathcal{X}/\mathcal{Y}} \to \mathcal{X}$
is separated,
\item $\Delta_{f, 1} = \Delta_f :
\mathcal{X} \to \mathcal{X} \times_\mathcal{Y} \mathcal{X}$
is separated, and
\item $\Delta_{f, 2} = e :
\mathcal{X} \to \mathcal{I}_{\mathcal{X}/\mathcal{Y}}$
is a closed immersion.
\end{enumerate}
\item
The following are equivalent
\begin{enumerate}
\item $\mathcal{I}_{\mathcal{X}/\mathcal{Y}} \to \mathcal{X}$
is quasi-separated,
\item $\Delta_{f, 1} = \Delta_f :
\mathcal{X} \to \mathcal{X} \times_\mathcal{Y} \mathcal{X}$
is quasi-separated, and
\item $\Delta_{f, 2} = e :
\mathcal{X} \to \mathcal{I}_{\mathcal{X}/\mathcal{Y}}$
is quasi-compact.
\end{enumerate}
\item
The following are equivalent
\begin{enumerate}
\item $\mathcal{I}_{\mathcal{X}/\mathcal{Y}} \to \mathcal{X}$
is locally separated,
\item $\Delta_{f, 1} = \Delta_f :
\mathcal{X} \to \mathcal{X} \times_\mathcal{Y} \mathcal{X}$
is locally separated, and
\item $\Delta_{f, 2} = e :
\mathcal{X} \to \mathcal{I}_{\mathcal{X}/\mathcal{Y}}$
is an immersion.
\end{enumerate}
\item
The following are equivalent
\begin{enumerate}
\item $\mathcal{I}_{\mathcal{X}/\mathcal{Y}} \to \mathcal{X}$
is unramified,
\item $f$ is DM.
\end{enumerate}
\item
The following are equivalent
\begin{enumerate}
\item $\mathcal{I}_{\mathcal{X}/\mathcal{Y}} \to \mathcal{X}$
is locally quasi-finite,
\item $f$ is quasi-DM.
\end{enumerate}
\end{enumerate}
\end{lemma}

\begin{proof}
Proof of (1), (2), and (3).
Choose an algebraic space $U$ and a surjective smooth morphism
$U \to \mathcal{X}$. Then
$G = U \times_\mathcal{X} \mathcal{I}_{\mathcal{X}/\mathcal{Y}}$
is an algebraic space over $U$ (Lemma \ref{lemma-inertia}).
In fact, $G$ is a group algebraic space over $U$
by the group law on relative
inertia constructed in Remark \ref{remark-inertia-is-group-in-spaces}.
Moreover, $G \to \mathcal{I}_{\mathcal{X}/\mathcal{Y}}$
is surjective and smooth as a base change of $U \to \mathcal{X}$.
Finally, the base change of
$e : \mathcal{X} \to \mathcal{I}_{\mathcal{X}/\mathcal{Y}}$
by $G \to \mathcal{I}_{\mathcal{X}/\mathcal{Y}}$
is the identity $U \to G$ of $G/U$.
Thus the equivalence of (a) and (c) follows from
Groupoids in Spaces, Lemma
\ref{spaces-groupoids-lemma-group-scheme-separated}.
Since $\Delta_{f, 2}$ is the diagonal of $\Delta_f$ we have
(b) $\Leftrightarrow$ (c) by definition.

\medskip\noindent
Proof of (4) and (5). Recall that (4)(b) means $\Delta_f$ is
unramified and (5)(b) means that $\Delta_f$ is locally quasi-finite.
Choose a scheme $Z$ and a morphism
$a : Z \to \mathcal{X} \times_\mathcal{Y} \mathcal{X}$.
Then $a = (x_1, x_2, \alpha)$ where $x_i : Z \to \mathcal{X}$
and $\alpha : f \circ x_1  \to f \circ x_2$ is a $2$-morphism.
Recall that
$$
\vcenter{
\xymatrix{
\mathit{Isom}_{\mathcal{X}/\mathcal{Y}}^\alpha(x_1, x_2)
\ar[d] \ar[r] &
Z \ar[d] \\
\mathcal{X} \ar[r]^{\Delta_f} &
\mathcal{X} \times_\mathcal{Y} \mathcal{X}
}
}
\quad\text{and}\quad
\vcenter{
\xymatrix{
\mathit{Isom}_{\mathcal{X}/\mathcal{Y}}(x_2, x_2)
\ar[d] \ar[r] &
Z \ar[d]^{x_2} \\
\mathcal{I}_{\mathcal{X}/\mathcal{Y}} \ar[r] &
\mathcal{X}
}
}
$$
are cartesian squares. By Lemma \ref{lemma-isom-pseudo-torsor-aut-over-space}
the
algebraic space $\mathit{Isom}_{\mathcal{X}/\mathcal{Y}}^\alpha(x_1, x_2)$
is a pseudo torsor for $\mathit{Isom}_{\mathcal{X}/\mathcal{Y}}(x_2, x_2)$
over $Z$. Thus the equivalences in (4) and (5) follow from
Groupoids in Spaces, Lemma
\ref{spaces-groupoids-lemma-pseudo-torsor-implications}.
\end{proof}

\end{document}
